\section{Introduction}
\label{sec:introduction}

Angular and kinked interfaces between dissimilar materials occur in many
adhesive and multimaterial joints. In such configurations, the interface
geometry itself can create local stress concentrations that strongly affect
damage initiation~\cite{GacoinEtAl2009}. In an idealized elastic setting, a
bimaterial interface corner is characterized by a power-law eigenfield whose
exponent and angular structure depend on the geometry and elastic contrast.
Classical existence conditions and characteristics of such singularities were
established for bonded isotropic wedges~\cite{BogyWang1971,DempseySinclair1981};
subsequent studies showed that changing the angle can substantially alter both
the singularity exponent and the intensity of the local field
\cite{MohammedLiechti2001}. Methods for determining singular-field parameters
in more complex three-dimensional configurations continue to be developed
\cite{MiyazakiFukudaNoda2025}, while small-scale relations between the
parameters of an undamaged corner field and those of a short crack emanating
from an interface corner have also been proposed
\cite{LyuZhangMai2026}. Thus, fracture problems depend not only on the
existence of a corner singularity itself, but also on how that singularity
interacts with a local nonlinear or process zone.

For initially undamaged bimaterial corners, the transition from the elastic
singularity to local fracture has been studied both through parameters of the
asymptotic field and through models with a finite cohesive or process zone.
Mohammed and Liechti combined experimental measurements with cohesive-zone
modeling of crack nucleation for different bimaterial-joint angles
\cite{MohammedLiechti2000}, whereas Labossiere and Dunn examined fracture
initiation at three-dimensional bimaterial corners
\cite{LabossiereDunn2001}. Akisanya and Meng showed that the use of a
critical parameter of the elastic corner field requires sufficient scale
separation between the nonlinear zone and the region dominated by the
asymptotic field~\cite{AkisanyaMeng2003}. In modern cohesive formulations,
the same requirement is expressed through the ratio of a characteristic
fracture-zone length to the geometric scale of the problem
\cite{GormanThouless2022,ThoulessGoutianos2023}. For local analytical models,
this directly motivates the small-scale assumption: the process zone should
be small compared with the scale of the outer singular field to which it is
matched.

A second group of studies concerns existing interface cracks and their
possible departure from the interface. For an ideally brittle description,
classical energetic analyses established conditions for initial crack kinking
out of an interface~\cite{HeHutchinson1989}; corresponding relations for
stress intensity factors and energy release rates were obtained for
anisotropic bimaterials~\cite{Wang1994}. Cohesive models show that path
selection may depend not only on fracture toughness, but also on cohesive
strength and the characteristic fracture-zone scale
\cite{ParmigianiThouless2006}. Analytical models of small process zones near
interface cracks provide the zone length, opening, and energy characteristics
\cite{KaminskyEtAl2023,KaminskyDudykChornoivan2025}; accounting for the
nonsingular $T$-stress further shows that higher-order terms of the outer
field can produce appreciable quantitative corrections
\cite{KaminskyDudykChornoivan2026}. These results provide the methodological
basis for the present process-zone approach, although the initial geometric
configuration considered here is different.

For corner points of kinked interfaces specifically, analytical models of
local fracture have been developed in several closely related settings.
Dudyk and Dikhtyarenko studied the initial stage of an interface crack
kinking from a corner point~\cite{DudykDikhtyarenko2012}, whereas Kaminsky,
Kipnis, and Polishchuk considered an initial fracture process zone arising
directly at the corner of a V-shaped interface
\cite{KaminskyKipnisPolishchuk2012}. For an interface crack with contacting
faces emanating from a corner point of a polygonal interface, parameters of a
small fracture zone were determined in
\cite{KaminskyDudykFenkiv2023}. A related formulation is the model of crack
initiation near a kinked interface in the presence of a plastic connecting
layer~\cite{KaminskyDudykPolishchuk2024}. Thus, the analytical framework for
local process zones near interface corners and interface cracks is well
developed; its application to the specific field generated by a stationary
contacting interfacial shear crack at the corner of a V-shaped interface
remains to be addressed.

Such an outer field was obtained by Nazarenko and Kipnis for a symmetric
interfacial shear crack with faces in frictionless contact
\cite{NazarenkoKipnis2022Stress}. In the present work, the two symmetric
branches of this defect are treated as a single V-shaped interfacial shear
crack. Each branch has length $l$, so the total length of the two branches is
$2l$. In the cracked configuration, a power-law singular field remains near
the corner; its exponent and amplitude depend on the geometry, elastic
contrast, and $l$. This field provides the immediate outer asymptotics for
the local process-zone problem considered below.

For the same configuration, Nazarenko and Kipnis also determined the stress
intensity factor for in-plane shear (Mode~II) at the outer tips of the
V-shaped crack, formulated its limiting-equilibrium condition, and showed
that this equilibrium is unstable with respect to further interfacial
propagation~\cite{NazarenkoKipnis2022Equilibrium}. That mechanism is localized
at the outer tips and is potentially competitive with the mechanism studied
here at the corner. In the present work, the V-shaped interfacial crack is
treated as a prescribed stationary configuration, and the applied loading is
assumed insufficient to cause loss of its local stability. Under these
conditions, the singular corner field can support a small tensile process
zone in one of the materials. The corner and outer-tip mechanisms are
therefore treated as competing mechanisms rather than as sequential stages
of a prescribed fracture scenario.

The subject of this work is the local process zone near the corner of a
V-shaped interfacial shear crack and the use of its characteristics to
predict branching initiation into one of the materials. Under the small-scale
approximation, the zone length in material $i$ ($i=1,2$) is denoted by $d_i$
and is assumed to be much smaller than the crack-branch length $l$. The
crack field without the zone is used as the outer field, and the leading-order
feedback of the zone on the outer crack tips is neglected. The objectives are
to establish in which material such a zone is physically admissible, determine
its equilibrium length $d_i$, opening $\delta_i$, work $W_i$ done against
the separation resistance, zone-related energy parameter
$J_i=\dd W_i/\dd d_i$, and the corresponding change in the local corner
singularity.

The aim is to construct an analytical model for the quasi-static family of
local equilibrium states of the process zone and, on that basis, formulate
conditions for branching initiation of the V-shaped interfacial shear crack
at its corner. Under plane strain and symmetric loading, a small tensile
process zone is considered in one of the materials and is oriented along the
bisector of the corresponding angular sector. The zone is modeled as a line
of normal displacement discontinuity subjected to a constant normal cohesive
traction. Mellin transformation, asymptotic matching, and the Wiener--Hopf
method are used to determine the conditions of physical admissibility, the
equilibrium zone length and opening, the work done against separation
resistance, the zone-related energy parameter, and the exponent of the
residual corner singularity. The opening and energy parameter may then be
used in deformation and energy criteria, respectively, for initiation of the
branched crack once the corresponding critical material parameter is
specified independently. Subsequent propagation and the finite branching path
are outside the scope of this local formulation.
